\documentclass[11pt]{article}
\usepackage[margin=1in]{geometry}
\usepackage[utf8]{inputenc}
\usepackage{microtype}
\usepackage{parskip}
\usepackage{amssymb}
\usepackage{hyperref}
\usepackage{shared/response}

\hypersetup{colorlinks=true, urlcolor=rrborder, linkcolor=black}

%% Custom bullet markers
\renewcommand{\labelitemi}{\textcolor{rrheader}{$\blacktriangleright$}}

\pagestyle{empty}

\begin{document}

%% ── Header block ──────────────────────────────────────────────────────────
\noindent
{\sffamily\bfseries\large\color{rrheader}
  scCCVGBen --- a single-cell centroid-coupled VAE benchmark}

\smallskip
\noindent
{\sffamily Zeyu Fu and co-authors \quad\textperiodcentered\quad
  School of Computer Science \quad\textperiodcentered\quad \today}

\vspace{4pt}
\noindent\textcolor{rrheader}{\rule{\linewidth}{0.6pt}}
\vspace{10pt}

%% ── Salutation ────────────────────────────────────────────────────────────
Dear Editor,

We submit a revised version of the manuscript previously titled
``CCVGAE: a centroid-coupled variational graph attention autoencoder for
stable and interpretable single-cell representation learning'' for further
consideration.

\medskip

{\itshape
The revision reflects a title change to \textbf{scCCVGBen} --- a single-cell
centroid-coupled VAE benchmark. The rename addresses three structural changes
prompted by reviewer feedback. The dataset cohort grew from 170 to 200
curated profiles (100 scRNA-seq, 100 scATAC-seq), giving the comparison axes
the statistical room they require. The encoder backbone, graph construction,
dataset, and evaluation suite were separated into independent axes so that
sensitivity to each factor can be read off independently. Figure provenance
was locked to the released benchmark manifest so that every panel can be
regenerated end-to-end without re-running training. The algorithmic
contribution is unchanged, but the manuscript now reads as a benchmark
companion rather than a method-introduction paper.}

\medskip

The principal architectural and experimental changes are as follows:

\begin{itemize}
  \item 200-dataset cohort with a full provenance manifest and public-safe
        supplementary table.
  \item Encoder-axis robustness across 14 backbones (Fig.~5).
  \item Graph-construction-axis robustness across 5 similarity definitions
        (Fig.~6).
  \item Cross-method scRNA benchmark with 12 completed external baselines
        (Fig.~7).
  \item Cross-method scATAC benchmark with 3 modality-specific baselines
        (Fig.~8).
  \item Three biological case studies with paired conditions (Figs.~9--11).
  \item Completed clustering-method sensitivity and seed-control clarifications
        addressing R1.2, R1.3 and R2.6 without relying on unreported entries.
  \item Hyperparameter-default justification and a documented sensitivity grid
        addressing R2.2 while keeping claims limited to completed benchmark
        results.
  \item Discussion of current foundation-model comparators, including scGPT and
        scFoundation, without adding unverified figure columns.
  \item Per-dataset summary table (Supplementary Table~S1) built from the
        released benchmark manifest --- addresses R1.1.
  \item Loss equation correction --- the adjacency-reconstruction term
        $w_{\mathrm{adj}} \cdot \|\hat{A} - A\|^2_F$ is now explicit in
        Equation~11 (it appeared in the architecture caption and Table~1 but
        was absent from the equation), addressing R2.1.
\end{itemize}

Two point-by-point rebuttal letters follow this cover letter, one per reviewer.

We thank the reviewers for their constructive comments, which materially
improved the scope and presentation of the work.

\vspace{16pt}
\noindent\textcolor{rrheader}{\rule{\linewidth}{0.4pt}}
\vspace{6pt}

\noindent
{\sffamily\bfseries Zeyu Fu} \quad (corresponding author)\\
{\small School of Computer Science\\
\href{mailto:peterponyusmith@gmail.com}{peterponyusmith@gmail.com}}

\vspace{4pt}
\noindent{\small Manuscript ID: scCCVGBen-2026 \quad\textperiodcentered\quad \today}

\end{document}
